\documentclass[a4paper]{article}

\usepackage{graphicx}
\usepackage{amsmath,amssymb}
\usepackage{minted}
\usepackage{multirow}
\usepackage{float}
\usepackage{todonotes}
\usepackage{forest}
\usepackage{xstring}
\usepackage{xcolor}
\usepackage[most]{tcolorbox}
\usepackage{xparse}
\usepackage{natbib}

\usetikzlibrary{graphs,graphdrawing}
\usegdlibrary{layered}

% for external graphviz
\input{/home/donkarlo/Dropbox/repo/nd_ai_project/src/nd_ai/knoweldge_representation/language/formal/computer/programming/tex/latex/code_clips/external_graphviz.tex}

% for tags
\input{/home/donkarlo/Dropbox/repo/nd_ai_project/src/nd_ai/knoweldge_representation/language/formal/computer/programming/tex/latex/parts/control_sequence/kind/semtag/semtag.tex}

% for links
\input{/home/donkarlo/Dropbox/repo/nd_ai_project/src/nd_ai/knoweldge_representation/language/formal/computer/programming/tex/latex/code_clips/seehere.tex}

\title{Saussure}
\date{}

\begin{document}
    \maketitle
    
    Saussure’s view \cite{saussure-1966-course-in-general-linguistics} is that language is not a list of names for things, but a system of signs, where each sign is formed by the relation between a signifier and a signified.
    
    The meaning or value of a sign does not come from itself alone, but from its differences and relations with other signs inside the whole language system.
    
    \begin{itemize}
        \item \textbf{Syntagmatic axis:} language elements combine in a linear sequence, for example words forming a sentence.
        \item \textbf{Paradigmatic axis:} language elements can replace each other in the same position, changing the meaning or value of the expression.
    \end{itemize}
    
    \bibliographystyle{plainnat}
    \bibliography{/home/donkarlo/Dropbox/repo/nd_shared_research/bibliography/bibliography.bib}
\end{document}